\documentclass{article}

\usepackage{neurips_2026}

\usepackage[utf8]{inputenc}
\usepackage[T1]{fontenc}
\usepackage{hyperref}
\usepackage{url}
\usepackage{booktabs}
\usepackage{amsmath}
\usepackage{amsfonts}
\usepackage{nicefrac}
\usepackage{microtype}
\usepackage{xcolor}
\usepackage{graphicx}
\usepackage[table]{xcolor}
\definecolor{cgreen}{HTML}{2ea043}
\definecolor{corange}{HTML}{d29922}
\definecolor{cred}{HTML}{cf222e}

\newcommand{\comment}[1]{\textcolor{orange}{[#1]}}
\def\eg{{\it e.g.}}
% shortcuts for technical terms
\def\R{\mathbb{R}}
\def\vrest{V^{\mathrm{rest}}}
\def\relu{\mathrm{ReLU}}

\title{Graph neural networks recover interpretable \comment{quantitative} circuit models from neural activity}

\author{%
  First Author \\
  Department / Institution \\
  \texttt{first@example.com} \\
  % \And
  % Second Author \\
  % Department / Institution \\
  % \texttt{second@example.com} \\
}

\begin{document}


\maketitle


\begin{abstract}

Synapse-level connectomes describe the structure of circuits, but not the electrical and chemical dynamics.  Conversely, large-scale recordings of neural activity capture these dynamics, but not the circuit structure.  We asked whether combining binary connectivity and recorded neural activity can be used to infer mechanistic models of neural circuits.  We trained a graph neural network model (GNN) to forecast the activity of \emph{Drosophila} visual system simulations.  Trained on activity trajectories in response to visual inputs, the model recovers effective connectivity weights, neuron types, and nonlinear activation functions, even when 400\% random connections are added to the adjacency matrix.  Moreover, it correctly predicts causal effects of connection removal, demonstrating the ability to infer mechanistic dependencies directly from activity data.

\end{abstract}


\section{Introduction}
\begin{itemize}
\item Learning dynamics from neural data
\item build on previous GNN framework tested on case studies, not biologically realistic 
\item Flyvis and Flywire more realistic (sparse, Dale law, structured, large, ...)
\item landscape of simulations -> many conditions, solutions, parameters -> agent in the loop HPO
\item benchmark shows good activity prediction but is not learning mechanistic model
\end{itemize}


\section{Related work}


Prior studies were successful in retrieving functional connectivity, functional properties of neurons, or predictive models of neural activity at the population level, but a method to infer an interpretable mechanistic model of complex neural assemblies is yet missing. \citet{mi_connectome-constrained_2021} infer the hidden voltage dynamics and functional connectivity from calcium recordings of parts of the \emph{C. elegans} nervous system using a simple biophysical neuron model.  While its predictive performance is modest, it demonstrates that integrating real activity recordings with connectivity can yield experimentally testable predictions beyond simulation.  \citet{pospisil_fly_2024} used the connectome of \emph{Drosophila melanogaster} \citep{dorkenwald_neuronal_2024} to create simulations of full brain activity and inferred effective connectivity from known perturbations of individual neurons with a simple linear model.  Lacking information about connectivity, \citet{wang_foundation_2025} trained a foundation model of mouse visual cortex activity on real two-photon recordings from 135,000 neurons, which predicts responses to both natural videos and out-of-distribution stimuli across animals.
\citet{prasse_predicting_2022} focus on accurate prediction without recovering connectivity while \citet{yu_discovering_2025} infer symbolic dynamical equations without validating connectivity weights. Both work with simulated data and do not learn neuronal heterogeneity.

Our method learns the functional properties of individual neurons and their relative connectivity weights, enabling transfer of the learned dynamics to different network configurations with modified connectivity or different neuron type compositions.

\subsubsection*{Simulations}
We simulated neural activity in the \emph{Drosophila} visual system using \textit{flyvis}' pretrained models \cite{lappalainen_connectome-constrained_2024}. The recurrent neural network contains $13{,}741$ neurons from $65$ cell types and $434{,}122$ synaptic connections. We restricted the original $721$ retinotopic columns to the central subset of $217$. Each neuron is modeled as a non-spiking compartment governed by
\begin{equation}
\tau_i\frac{dv_i(t)}{dt} = -v_i(t) + V_i^{\text{rest}}
+ \sum_{j\in\mathcal{N}_i}\mathbf{W}_{ij}\,\text{ReLU}\!\big(v_j(t)\big)
+ I_i(t) + \sigma\,\xi_i(t),
\label{eq:simulated_update}
\end{equation}
where $\tau_i$ and $V_i^{\mathrm{rest}}$ are cell-type parameters, $\mathbf{W}_{ij}$ is the connectome-constrained synaptic weight, $I_i(t)$ the visual input, and $\xi_i(t)\sim\mathcal{N}(0,1)$ is independent Gaussian noise scaled by $\sigma$. This intrinsic noise perturbs the dynamical trajectory of $v_i(t)$ as it propagates through the network. We also added measurement noise post hoc, independent Gaussian noise scaled by $\gamma$ We generated data at three intrinsic noise levels and with two measuremnt noise levels : $\sigma\in\{0,\,0.05,\,0.5\}, \gamma \in\{0,\,0.1\}$. 

\comment{Janne complete with Flywirte model}

Visual inputs are derived from the DAVIS dataset \cite{davis1,davis2} \comment{Vijay add YouTube-VOS, discuss training/test/validation}, natural video sequences center-cropped to $60\%$, converted to grayscale, and resampled onto the hexagonal photoreceptor lattice of $217$ columns via a Gaussian box-eye filter, giving $1{,}736$ input neurons ($8$ photoreceptor types R1--R8 per column). Sequences are augmented with horizontal/vertical flips and four rotations ($0°, 90°, 180°, 270°$), and split $80/20$ into train/validation without leakage. 



\subsubsection*{Graph neural network model}
We approximate the simulated voltage dynamics by a message-passing GNN \cite{gilmer_neural_2017}
\begin{equation}
\frac{d\hat{v}_i(t)}{dt}
= f_\theta\!\Big(v_i(t),\,\mathbf{a}_i,\,
\sum_{j\in\mathcal{N}_i} \widehat{\mathbf{W}}_{ij}\,
g_\phi\!\big(v_j(t),\mathbf{a}_j\big)^2,\,
I_i(t)\Big).
\label{eq:gnn_ct}
\end{equation}
Nodes of the GNN correspond to neurons and edges correspond to functional synaptic connections. The GNN learns a latent embedding $\mathbf{a}_i\in\mathbb{R}^2$ for each neuron $i$, giving each neuron a compact latent identity to capture cell-type-specific properties (like time constants and nonlinearities). Neuron update $f_\theta=\text{MLP}_0$ and edge message $g_\phi=\text{MLP}_1$ are three-layer perceptrons (width $64$, ReLU, linear output). $g_\phi$ maps presynaptic inputs $v_j$ to nonnegative messages (via squaring) depending on neural embedding $\mathbf{a}_j$, weighted by $\widehat{\mathbf{W}}_{ij}$. $f_\theta$ processes the postsynaptic voltage $v_i$, aggregated input, and external input $I_i(t)$ to predict $d\hat{v}_i(t)/dt$, depending on neural embedding $\mathbf{a}_i$. During training, inputs $I_i(t)$, adjacency $\mathcal{N}_i$, and activity $v_i(t)$ are given. The MLPs, $\widehat{\mathbf{W}}_{ij}$, and $\mathbf{a}_i$ are optimized by minimizing
\begin{equation}
\mathcal{L}_{\mathrm{pred}} = \sum_{i,t}\|\hat{y}_i(t) - y_i(t)\|_2
\label{eq:loss_pred}
\end{equation}
between simulator $y_i(t)=dv_i(t)/dt$ and GNN predictions $\hat{y}_i(t)=d\hat{v}_i(t)/dt$.

\subsubsection*{Regularization}
Many synaptic weight configurations give rise to identical neural activity, making the recovery of $\widehat{\mathbf{W}}$ from voltage traces alone fundamentally non-unique (see Appendix \ref{app:degeneracy}). To constrain the solution space by adding regularization terms to $\mathcal{L}_{\mathrm{pred}}$ 
\begin{eqnarray}
\mathcal{L} &=& \mathcal{L}_{\mathrm{pred}}
+ \lambda_0\|\theta\|_1 + \lambda_1\|\phi\|_1
+ \gamma_0\|\widehat{\mathbf{W}}\|_1 + \gamma_1\|\widehat{\mathbf{W}}\|_2 \nonumber\\
&& +\; \mu_0\!\left\|\mathrm{ReLU}\!\left(-\tfrac{\partial\,g_\phi(v,\mathbf{a})}{\partial v}\right)\right\|_2
+ \mu_1\|g_\phi(v_\star,\mathbf{a})-v_\star\|_2.
\label{eq:loss_full}
\end{eqnarray}
The $\ell_1$ penalties on $\theta$ and $\phi$ promote sparsity in the MLP weights, biasing $f_\theta$ and $g_\phi$ toward simpler functions. The $\ell_2$ penalties prevent weight explosion and stabilize training. The $\ell_1$ term $\gamma_0$ promotes sparse effective connectivity, while the $\ell_2$ term $\gamma_1$ penalizes large synaptic weights. The monotonicity term penalizes negative slopes of $g_\phi$ with respect to $v$, encoding the prior that synaptic drive increases with presynaptic voltage. The normalization term anchors $g_\phi$ at a reference voltage $v_\star$, removing the scale degeneracy between $g_\phi$ and $\widehat{\mathbf{W}}$. The coefficients $\lambda_0,\lambda_1,\gamma_0,\gamma_1,\mu_0,\mu_1$ are hyperparameters tuned via the agentic optimization procedure described below.

\subsubsection*{Agent in the loop parameters optimization}
The combined space of simulations and machine learning model parameters is too large to be explored with a grid search. Instead we deployed a closed-loop system 
where a Claude Code agent interpreted experiment results, maintained a structured summary, and proposed next intervention on the simulation and/or machine learning parameters (see Appendix \ref{app:agent}). The primary optimization target is either identifiability of the connectivity weights or rollout prediction accuracy. The agent is instructed to perform scientific reasoning: testable hypotheses must be drawn, repeatable experiments are conducted to validate or falsify them. The agent is told not to predict the outcome of a model training. Only training results can validate or falsify its hypotheses. Beyond  optimal parameter configurations, the agent can be queried at any point to report its accumulated understanding, e.g. established principles, causal chains, and open questions that persist.


\subsubsection*{Benchmark}

In the Known-ODE baseline, both the ordinary differential equation (ODE) and the ReLU activation are fixed to the ground-truth form (Eq.~\ref{eq:simulated_update}). We test the ability to learn per-neuron time constants $\hat{\tau}_i$, resting potentials $\hat{V}_i^{\mathrm{rest}}$, and synaptic weights $\widehat{W}_{ij}$ if the model is given.





\section{Results}

\begin{itemize}
\item noise free, noise 0.05, noise 0.5
\item rollout with without noise
\item ablation edges -> mechanistic test
\item parameter extraction - embedding type given \textbf{NEW}
\item unique config (400K extra edges) test on all conditions
\item ablation test
\item joint learning GNN + stimuli
\item extra edges +1.6M
\item missing edges 20\%

\paragraph{Robustness to erroneous adjacency matrices.}
In practice, annotated connectome reconstructions contain both false negatives (missing
connections) and false positives (spurious connections). To test robustness to
both error types, we \textit{(i)} augmented the ground-truth adjacency matrix
with randomly added null edges ($+400\%$; $434{,}112 \to 2{,}170{,}560$ edges)
and \textit{(ii)} randomly removed edges ($-20\%$; $434{,}112 \to 347{,}000$
edges). In both cases the GNN was trained on the corrupted connectivity while
neural activity was simulated with the full ground-truth connectome
($434{,}112$ edges).



\item measurement noise \textbf{NEW}
\item partial data 1/5 frames only \textbf{NEW}
\item partial data 10\% neurons removed \textbf{NEW}
\item benchmark Known-ODE MLP EED  \textbf{NEW}
\end{itemize}

\newcommand{\good}[1]{\textcolor{green!50!black}{#1}}


\begin{table}[h]
\centering
\caption{The GNN evaluation results are given over five-fold cross-validation
(mean~$\pm$~SD, $n{=}5$ seeds). The GNN training parameters have been specifically tuned by an agentic loop for each condition. Prediction: DAVIS-trained model tested on held-out YouTube-VOS sequences; one-step and rollout Pearson~$r$ over 8\,000 noise-free test frames.
Parameter recovery: model re-trained with YouTube-VOS sequences; $R^2_{\widehat{W}}$ over all edges, $R^2_{\tau}$ and $R^2_{V^{\mathrm{rest}}}$ over all neurons,
GMM clustering accuracy over $65$ cell types. 
\textcolor{green!50!black}{Green}: value $>0.9$.}
\label{tab:cv_per_condition}
\footnotesize
\setlength{\tabcolsep}{4pt}
\begin{tabular}{lccrrrrrr}
\toprule
& & & \multicolumn{2}{c}{prediction} & \multicolumn{4}{c}{parameter recovery} \\
\cmidrule(lr){4-5}\cmidrule(lr){6-9}
condition & noise $\sigma$ & edges
  & one-step $r$ & rollout $r$
  & $R^2_{\widehat{W}}$ & $R^2_{\tau}$ & $R^2_{V^{\mathrm{rest}}}$ & cluster\ acc. \\
\midrule
noise-free   & $0$    & 434\,112
  & \good{$0.99{\pm}0.00$} & \good{$0.99{\pm}0.00$}
  & \good{$0.92{\pm}0.01$} & $0.43{\pm}0.34$ & $0.07{\pm}0.02$ & $0.78{\pm}0.02$ \\
low noise& $0.05$ & 434\,112
  & \good{$0.98{\pm}0.00$} & \good{$0.98{\pm}0.00$}
  & \good{$0.96{\pm}0.04$} & \good{$0.99{\pm}0.00$} & $0.44{\pm}0.18$ & $0.84{\pm}0.03$ \\
high noise & $0.5$  & 434\,112
  & \good{$0.99{\pm}0.00$} & \good{$0.99{\pm}0.00$}
  & \good{$1.00{\pm}0.00$} & \good{$1.00{\pm}0.00$} & \good{$0.98{\pm}0.01$} & \good{$0.90{\pm}0.01$} \\
$+400\%$ null edges & $0.05$ & 2\,170\,560
  & \good{$0.98{\pm}0.00$} & \good{$0.98{\pm}0.00$}
  & \good{$0.97{\pm}0.01$} & $0.78{\pm}0.39$ & $0.02{\pm}0.01$ & $0.57{\pm}0.04$ \\[2pt]
$-20\%$ edges removed & $0.05$ & $347\,000$
  & \good{$0.99{\pm}0.01$} & \good{$0.99{\pm}0.01$}
  & $0.83{\pm}0.07$ & \good{$0.91{\pm}0.10$} & $0.39{\pm}0.11$ & $0.55{\pm}0.04$ \\
\bottomrule
\end{tabular}
\end{table}

\newcommand{\good}[1]{\textcolor{green!50!black}{#1}}
\newcommand{\comment}[1]{\textcolor{gray}{\textit{#1}}}

\begin{table}[h]
\centering
\caption{The GNN evaluation results are given over five-fold cross-validation
(mean~$\pm$~SD, $n{=}5$ seeds). The model was trained with parameters obtained by the agentic loop with condition $400\%$ null-edges (Table \ref{tab:cv_per_condition}). Prediction: DAVIS-trained model tested on held-out YouTube-VOS sequences; one-step and rollout Pearson~$r$ over 8\,000 noise-free test frames.
Parameter recovery: model re-trained with YouTube-VOS sequences; $R^2_{\widehat{W}}$ over all edges, $R^2_{\tau}$ and $R^2_{V^{\mathrm{rest}}}$ over all neurons,
GMM clustering accuracy over $65$ cell types. 
\textcolor{green!50!black}{Green}: value $>0.9$.}
\label{tab:cv_cross_noise}

\footnotesize
\setlength{\tabcolsep}{4pt}
\begin{tabular}{lccrrrrrr}
\toprule
& & & \multicolumn{2}{c}{prediction} & \multicolumn{4}{c}{parameter recovery} \\
\cmidrule(lr){4-5}\cmidrule(lr){6-9}
condition & noise $\sigma$ & edges
  & one-step $r$ & rollout $r$
  & $R^2_{\widehat{W}}$ & $R^2_{\tau}$ & $R^2_{V^{\mathrm{rest}}}$ & cluster.\ acc. \\
\midrule
noise-free          & $0$    & $434\,112$
  & \good{$0.99{\pm}0.00$} & $0.66{\pm}0.40$
  & $0.86{\pm}0.02$ & $0.05{\pm}0.04$ & $0.11{\pm}0.03$ & $0.73{\pm}0.04$ \\
low noise  & $0.05$ & $434\,112$
  & \good{$0.99{\pm}0.01$} & \good{$0.98{\pm}0.01$}
  & \good{$0.95{\pm}0.04$} & \good{$0.98{\pm}0.01$} & $0.42{\pm}0.12$ & $0.83{\pm}0.04$ \\
high noise  & $0.5$  & $434\,112$
  & \good{$0.99{\pm}0.01$} & \good{$0.99{\pm}0.01$}
  & \good{$0.99{\pm}0.01$} & \good{$1.00{\pm}0.00$} & $0.71{\pm}0.36$ & \good{$0.93{\pm}0.03$} \\
$+400\%$ null edges & $0.05$ & $2\,170\,560$
  & \good{$0.98{\pm}0.00$} & \good{$0.98{\pm}0.00$}
  & \good{$0.97{\pm}0.01$} & $0.78{\pm}0.39$ & $0.02{\pm}0.01$ & $0.57{\pm}0.04$ \\[2pt]
$-20\%$ edges removed & $0.05$ & $347\,000$
  & \good{$0.96{\pm}0.00$} & \good{$0.92{\pm}0.00$}
  & $0.85{\pm}0.00$ & \good{$0.94{\pm}0.01$} & $0.24{\pm}0.05$ & $0.53{\pm}0.01$ \\
\bottomrule
\end{tabular}
\end{table}


\begin{table}[h]
\centering
\caption{The joint GNN+INR evaluation results are given over five-fold validation (mean ± SD, n=5 seeds). The model was trained with parameters obtained by the agentic loop on low-noise condition (DAVIS sequences, Table \ref{tab:cv_per_condition}). Prediction: one-step and rollout Pearson $r$ over 8 000 noise-free frames. Stimuli $r$: Pearson $r$ between the INR output and the true stimuli. Parameter recovery: model re-trained on held-out video sequences; $R^2_{\widehat{W}}$ over all edges, $R^2_{\tau}$ and $R^2_{V^{\mathrm{rest}}}$ over all neurons,
GMM clustering accuracy over $65$ cell types.
\textcolor{green!50!black}{Green}: value $>0.9$.}
\label{tab:cv_inr}
\footnotesize
\setlength{\tabcolsep}{3pt}
\begin{tabular}{lccrrrrrrrrr}
\toprule
& & & \multicolumn{3}{c}{prediction} & \multicolumn{4}{c}{parameter recovery} \\
\cmidrule(lr){4-6}\cmidrule(lr){7-10}
condition & noise $\sigma$ & edges
  & one-step $r$ & rollout $r$ & stimuli $r$
  & $R^2_{\widehat{W}}$ & $R^2_{\tau}$ & $R^2_{V^{\mathrm{rest}}}$ & cluster\ acc. \\
\midrule
DAVIS & $0.05$ & $434\,112$
  & \comment{pending} & \comment{pending} & \comment{pending}
  & \comment{pending} & \comment{pending} & \comment{pending} & \comment{pending} \\
YouTube-VOS  & $0.05$ & $434\,112$
  & \good{$0.943{\pm}0.003$} & $0.42{\pm}0.13$ & $0.61{\pm}0.26$
  & $0.89{\pm}0.05$ & $0.84{\pm}0.09$ & $0.16{\pm}0.04$ & $0.78{\pm}0.05$ \\
\bottomrule
\end{tabular}
\end{table}



\begin{table}[h]
\centering
\caption{The Known-ODE evaluation results are given over five-fold cross-validation
(mean~$\pm$~SD, $n{=}5$ seeds). The Known-ODE training parameters have been specifically tuned by an agentic loop for each condition. Prediction: DAVIS-trained model tested on held-out YouTube-VOS sequences; one-step and rollout Pearson~$r$ over 8\,000 noise-free test frames.
Parameter recovery: model re-trained on YouTube-VOS sequences; $R^2_{\widehat{W}}$ over all edges, $R^2_{\tau}$ and $R^2_{V^{\mathrm{rest}}}$ over all neurons,
GMM clustering accuracy over $65$ cell types.
\textcolor{green!50!black}{Green}: value $>0.9$.}
\label{tab:cv_known_ode}
\footnotesize
\setlength{\tabcolsep}{4pt}
\begin{tabular}{lccrrrrrr}
\toprule
& & & \multicolumn{2}{c}{prediction} & \multicolumn{4}{c}{parameter recovery} \\
\cmidrule(lr){4-5}\cmidrule(lr){6-9}
condition & noise $\sigma$ & edges
  & one-step $r$ & rollout $r$
  & $R^2_{\widehat{W}}$ & $R^2_{\tau}$ & $R^2_{V^{\mathrm{rest}}}$ & cluster\ acc. \\
\midrule
noise-free                      & $0$    & $434\,112$
  & \good{$1.00{\pm}0.00$} & \good{$1.00{\pm}0.00$}
  & \good{$0.97{\pm}0.00$} & $0.004{\pm}0.000$ & \good{$0.92{\pm}0.00$} & $0.61{\pm}0.00$ \\
low noise             & $0.05$ & $434\,112$
  & \good{$1.00{\pm}0.00$} & \good{$1.00{\pm}0.00$}
  & \good{$0.99{\pm}0.00$} & \good{$1.00{\pm}0.00$} & \good{$0.97{\pm}0.00$} & $0.63{\pm}0.01$ \\
high noise                 & $0.5$  & $434\,112$
  & \good{$1.00{\pm}0.00$} & \good{$1.00{\pm}0.00$}
  & \good{$1.00{\pm}0.00$} & \good{$1.00{\pm}0.00$} & \good{$1.00{\pm}0.00$} & $0.61{\pm}0.01$ \\
intrinsic + meas. noise   & $0.05$ & $434\,112$
  & \good{$0.95{\pm}0.00$} & \good{$0.93{\pm}0.00$}
  & $0.79{\pm}0.00$ & $0.81{\pm}0.00$ & $0.18{\pm}0.00$ & $0.42{\pm}0.00$ \\[2pt]
$+400\%$ null edges             & $0.05$ & $2\,170\,560$
  & \good{$0.99{\pm}0.00$} & \good{$0.99{\pm}0.00$}
  & \good{$0.93{\pm}0.00$} & \good{$0.99{\pm}0.00$} & $0.72{\pm}0.00$ & $0.37{\pm}0.01$ \\
$-20\%$ edges removed           & $0.05$ & $347\,000$
  & \good{$0.99{\pm}0.00$} & \good{$0.98{\pm}0.00$}
  & \good{$0.90{\pm}0.00$} & \good{$0.99{\pm}0.00$} & $0.42{\pm}0.00$ & $0.55{\pm}0.03$ \\
\bottomrule
\end{tabular}
\end{table}





\section{Discussion}
\begin{itemize}
\item The present framework propose a set of neural activity data, ML models and standardized test to developp and assess learning neural dynamics from data
\item prediction activity together with parameters extraction is a necessary (standard) to assess dynamics learning: good prediction but not good model, not good prediction but good model
\item we stress that measurement noise, partial data are current bottlenecks of the approach
\end{itemize}

\section{Conclusion}

\begin{ack}
  % Remove this section for anonymous submission.
\end{ack}

% \section*{References}

{\small
\bibliographystyle{plainnat}
\bibliography{references}
}
\appendix
% appendix.tex — \input{appendix} before \end{document}
% appendix.tex — \input{appendix} before \end{document}

\section{Degeneracy of the inverse problem}
\label{app:degeneracy}
In the following, we characterise the dynamical equivalence class of $\mathbf{W}$, the set of all weight matrices that reproduce the observed trajectories. We show that Eq.~\eqref{eq:simulated_update} can be linearised per neuron, that this linear system admits a solution, and finally that it admits many different solutions.

\paragraph{Linearisation of the simulation equation.}
Rearranging Eq.~\eqref{eq:simulated_update} to isolate the incoming synaptic
weights gives
\begin{equation}
  \sum_{j\in\mathcal{N}_i} \mathbf{W}_{ij}\,\text{ReLU}\!\big(v_j(t)\big)
  = \tau_i \frac{dv_i(t)}{dt} + v_i(t) - V_i^{\text{rest}} - I_i(t).
\end{equation}
Defining the presynaptic activity $h_j(t) = \text{ReLU}(v_j(t))$ and the
residual $b_i(t) = \tau_i \frac{dv_i(t)}{dt} + v_i(t) - V_i^{\text{rest}} -
I_i(t)$, and stacking the $T$ time-point constraints into matrix form, yields
an independent linear system for each postsynaptic neuron $i$:
\begin{equation}
  \mathbf{H}_i\,\mathbf{w}_i = \mathbf{b}_i,
  \label{eq:linear_system}
\end{equation}
where $\mathbf{H}_i \in \mathbb{R}^{T \times d_i}$ is the presynaptic
activity matrix, $d_i$ is the in-degree of neuron $i$, $\mathbf{w}_i \in
\mathbb{R}^{d_i}$ the incoming weight vector, and $\mathbf{b}_i \in
\mathbb{R}^T$ the stacked residuals. This holds for each postsynaptic neuron
$i = 1, \ldots, N$, where $N = 13{,}697$ and $d_i$ ranges from $1$ to $208$.
The null space $\ker(\mathbf{H}_i) = \{\boldsymbol{\delta} \in \mathbb{R}^{d_i}
: \mathbf{H}_i\boldsymbol{\delta} = \mathbf{0}\}$ characterises the
degeneracy: any $\mathbf{w}_i + \boldsymbol{\delta}$ with $\boldsymbol{\delta}
\in \ker(\mathbf{H}_i)$ satisfies Eq.~\eqref{eq:linear_system} equally and
produces identical dynamics. In the following, we characterise the
dynamical equivalence class of $\mathbf{W}$, the set of all weight
matrices that reproduce the observed trajectories.


\paragraph{Pseudoinverse recovery.}
Assuming $\tau_i$ and $V_i^{\mathrm{rest}}$ are known, Eq.~\eqref{eq:linear_system}
admits the minimum-norm solution $\mathbf{w}_i^+ = \mathbf{H}_i^+\mathbf{b}_i$.
Table~\ref{tab:pseudoinverse} shows that different solvers achieve rollout
$R^2 = 1.000$ regardless of $R^2_{\mathbf{W}}$: many weight configurations
produce identical dynamics. Yet the solutions remain close to each other, 
$R^2_{\mathbf{W}} \geq 0.988$ for the truncated SVD. We also found that
$R^2_{\mathbf{W}}$ increases monotonically with $\sigma$; independent noise
realisations appear to shrink $\ker(\mathbf{H}_i)$.

\begin{table}[ht!]
\centering\small
\begin{tabular}{llrrr}
\toprule
Method & noise $\sigma$ & Conn $R^2_{\mathbf{W}}$ & Rollout $R^2$ & \% edges in $\ker(\mathbf{H}_i)$ \\
\midrule
Truncated SVD   & $0$    & 0.99 & 1.00 & 26.5\% \\
                & $0.05$ & 0.99 & 1.00 & 2.4\%  \\
                & $0.5$  & 1.00 & 1.00 & 0.0\%  \\
\midrule
Ridge ($\lambda=1$) & $0$    & 0.88 & 1.00 & 26.5\% \\
                    & $0.05$ & 0.99 & 1.00 & 2.4\%  \\
                    & $0.5$  & 1.00 & 1.00 & 0.0\%  \\
\midrule
CV-ridge        & $0$    & 0.99 & 1.00 & 26.5\% \\
                & $0.05$ & 0.99 & 1.00 & 2.4\%  \\
                & $0.5$  & 1.00 & 1.00 & 0.0\%  \\
\bottomrule
\end{tabular}
\caption{Weight recovery from Eq.~\eqref{eq:linear_system} with known
  $\tau_i$, $V_i^{\mathrm{rest}}$, using three solvers: truncated pseudoinverse
  (singular values $>10^{-12}$), ridge regression with fixed $\lambda=1$, and
  CV-ridge with per-neuron $\lambda$ selected by cross-validation.}
\label{tab:pseudoinverse}
\end{table}

\paragraph{Structural origin of the null space.}
In fact, the columnar organisation of the connectome directly implies the
existence of $\ker(\mathbf{H}_i)$, without SVD argument. The flyvis circuit
spans 217 retinotopic columns, and each postsynaptic neuron $i$ typically
receives input from $k_{i\alpha} > 1$ neurons of the same cell type $\alpha$
across columns. A moving stimulus crosses the eye so that column $c$ sees
what column $c-1$ saw moments earlier: same-type neurons produce time-shifted
copies of the same signal, giving the corresponding $k_{i\alpha}$ columns of
$\mathbf{H}_i$ rank $\approx 1$ instead of $k_{i\alpha}$. Any sum-zero
perturbation $\sum_{m=1}^{k_{i\alpha}} \delta_{j_m} = 0$ on those weights
leaves $\mathbf{H}_i \mathbf{w}_i$ unchanged, contributing $k_{i\alpha} - 1$
free directions to $\ker(\mathbf{H}_i)$. Summing over all neurons and cell
types:
\begin{equation}
  \dim\ker(\mathbf{H}) = \sum_{i=1}^{N}\sum_{\alpha:\,k_{i\alpha}>1}(k_{i\alpha}-1)
  \approx 121{,}100 \quad (\sim 28\%\ \text{of edges}).
  \label{eq:null_structural}
\end{equation}

\paragraph{Empirical confirmation of many solutions.}
To verify Eq.~\eqref{eq:null_structural} without any assumption, we
generated 960 perturbed connectivity matrices $\mathbf{W}' =
\mathbf{W} + \boldsymbol{\Delta}$, where $\boldsymbol{\Delta}$ applies
sum-preserving perturbations within same-type groups. 13 cell types lack duplicate-type groups, there are no suchvariants for theses cell types. For the other 51 cell-types, we established variants, computed $R^2_{\mathbf{W}}$ (weight
preservation) and rollout $R^2$ (ODE trajectory preservation over 1,000
frames). $R^2_{\mathbf{W}}$ ranges from 0.28 to 1.0 across variants while
rollout $R^2$ remains above 0.96 throughout. Perturbed weights differ substantially from the ground truth yet produce nearly identical dynamics (Table~\ref{tab:lambda_8}). Further, we looked for the minimal-edges $\mathbf{W}$ by collapsing each degenerate $(i,\alpha)$ group of $k_{i\alpha}$ edges onto a
single representative edge $e^* = \arg\max_e |W_e^*|$, zeroing the remaining
$k_{i\alpha}-1$ edges. We tested two variants: sum-preserving
($W_{e^*} = \sum_e W_e^*$, exact null-space member) and calibrated
($W_{e^*}$ minimising the per-group dynamics residual via one least-squares
scalar fit). Both zero $308{,}160$ edges ($71\%$ of $434{,}112$). The sum-preserving variant achieves rollout $R^2 = 0.94$ with $R^2_{\mathbf{W}} = 0.10$.
The calibrated variant achieves rollout $R^2 = 0.98$ with $R^2_{\mathbf{W}} = 0.39$.

\begin{table}[ht!]
\centering\footnotesize
\begin{tabular}{lrr@{\quad}lrr@{\quad}lrr}
\toprule
Cell Type & Conn $R^2_{\mathbf{W}}$ & Rollout $R^2$ && & & & & \\
\midrule
R2     & \cellcolor{corange}0.94 & 1.00 & CT1(M10) & \cellcolor{corange}0.88 & 1.00 & Tm4   & \cellcolor{cgreen}0.99 & 1.00 \\
R3     & \cellcolor{cgreen}0.97  & 1.00 & Mi2      & \cellcolor{corange}0.91 & 0.98 & Tm5Y  & \cellcolor{cgreen}1.00 & 1.00 \\
R4     & \cellcolor{corange}0.74 & 1.00 & Mi3      & \cellcolor{corange}0.94 & 1.00 & Tm5a  & \cellcolor{cgreen}0.97 & 1.00 \\
R5     & \cellcolor{cgreen}0.99  & 1.00 & Mi14     & \cellcolor{cgreen}0.96 & 1.00 & Tm5c  & \cellcolor{cred}0.28   & 1.00 \\
R6     & \cellcolor{cgreen}1.00  & 1.00 & Mi15     & \cellcolor{corange}0.87 & 1.00 & Tm16  & \cellcolor{cgreen}0.99 & 1.00 \\
R7     & \cellcolor{cgreen}1.00  & 1.00 & T2       & \cellcolor{cgreen}0.97  & 0.99 & Tm20  & \cellcolor{cgreen}1.00 & 1.00 \\
R8     & \cellcolor{corange}0.90 & 1.00 & T2a      & \cellcolor{cgreen}0.97  & 1.00 & Tm28  & \cellcolor{corange}0.75 & 0.99 \\
L1     & \cellcolor{corange}0.73 & 1.00 & T3       & \cellcolor{cgreen}0.97  & 1.00 & Tm30  & \cellcolor{corange}0.57 & 1.00 \\
L2     & \cellcolor{corange}0.62 & 1.00 & T4a      & \cellcolor{cgreen}0.97  & 1.00 & TmY3  & \cellcolor{cgreen}1.00 & 1.00 \\
L4     & \cellcolor{cgreen}1.00  & 1.00 & T4b      & \cellcolor{cgreen}0.99  & 1.00 & TmY4  & \cellcolor{cgreen}0.99 & 1.00 \\
L5     & \cellcolor{corange}0.73 & 0.99 & T4c      & \cellcolor{cgreen}0.99  & 1.00 & TmY5a & \cellcolor{corange}0.94 & 1.00 \\
Lawf1  & \cellcolor{cgreen}1.00  & 1.00 & T4d      & \cellcolor{cgreen}0.98  & 1.00 & TmY9  & \cellcolor{cgreen}0.95 & 1.00 \\
Am     & \cellcolor{cgreen}0.95  & 1.00 & T5a      & \cellcolor{cgreen}1.00  & 1.00 & TmY10 & \cellcolor{cgreen}0.97 & 0.96 \\
C2     & \cellcolor{corange}0.93 & 1.00 & T5b      & \cellcolor{cgreen}1.00  & 1.00 & TmY13 & \cellcolor{cgreen}0.99 & 1.00 \\
\bottomrule
\end{tabular}
\caption{Single-type variants for 51 cell types obtained. For each variant we
computed R2 W (weight preservation) and rollout R2 (ODE trajectory preservation over 1,000 frames). Green ($R^2>0.95$), orange ($0.50<R^2\leq 0.95$), red ($R^2\leq 0.50$).}
\label{tab:lambda_8}
\end{table}



\paragraph{Summary.}
Eq.~\eqref{eq:simulated_update} linearises into independent per-neuron systems
Eq.~\eqref{eq:linear_system}, each admitting many solutions as shown by three independent approaches. Pseudoinverse recovery, connectivity variants, and minimal-support collapse; all
show that dynamics can be preserved far from the ground-truth $\mathbf{W}$.





\section{Sensitivity analysis of the FlyVis ODE}


%%% Goal of this section: explain why the known ODE method does not recover tau %%%

Given an ODE, sensitivity analysis [TODO: right term?][TODO: ref] helps identify directions in parameter space that minimally impact the solution. The solution $x_i(t)$ to a general ODE
\begin{equation}
\dot{x_i} = f_i(x, \theta)    
\end{equation}
with the state $x \in \R^n$ and parameters $\theta \in \R^p$, implicitly depends on $\theta$. The \emph{sensitivity} $S_{i\alpha} := \partial x_i / \partial \theta_\alpha$ of the $i$-th component to the specific parameter direction $\theta_\alpha$ is the solution to another ODE:
\begin{equation}
\dot{S_{i; \alpha}} = \sum_j \frac{\partial f_i}{\partial x_j} S_{j; \alpha} + \frac{\partial f_i}{\partial \theta_\alpha} \label{eq:sens-ode}
\end{equation}
The ODE is initialized $x(t=0) = x_0$ to a value independent of the parameters $\theta$ and therefore we have $S_{i; \alpha}(t=0)= 0$. The second term on the right hand side describes the source term that causes $S_{i; \alpha}$ to gradually take on non-zero values. In the absence of the source term $S_{i; \alpha} = 0$ for all time. If we pick a direction in parameter space, defined by a unit vector $\lambda_\alpha$ we can describe the evolution of $s_i(t) := \sum_\alpha S_{i; \alpha} \lambda_\alpha$.

The noise-free FlyVis ODE \ref{eq:simulated_update} can be rewritten
\begin{equation}
    \dot{x_i} = -\omega_i x_i + U_i + \sum_j G_{ij} \relu(x_j) + \omega_i I_i
\end{equation}
with a suitable redefinition of the parameters (TODO: fix)
\begin{equation}
    \omega_i = 1/\tau_i, \ U_i = \vrest_i \tau_i, \ W_{ij} = G_{ij} \tau_i
\end{equation}
The Jacobian $J(x)$, which appears in \ref{eq:sens-ode}, has the special form $J(x) = \Omega + G\odot H(x)$ in matrix notation where $H(x) = d \relu (x) / dx$ and $\odot$ denotes elementwise multiplication. The Jacobian is a constant matrix in any region of state space where the signs of $x_i$ are unchanged for all $i$, \eg, $J(x) = \Omega$, for $x \in \R_-^n$. We compute the second term in \ref{eq:sens-ode} for each parameter below:
\begin{eqnarray}
    \omega_j & : & \frac{\partial f_i}{\partial \omega_j} = (-x_i + I_i) \delta_{ij} \\
    U_j & : & \frac{\partial f_i}{\partial U_j} = \delta_{ij} \\
    G_{jk} & : & \frac{\partial f_i}{\partial G_{jk}} = \relu(x_k) \delta_{ij}
\end{eqnarray}
Note that $I_i$ is non-zero only for the input neurons. The training data have the property that 
% If the jacobian is a negative semi-definite matrix for a majority of parameter space, The evolution of $S_{i; \alpha}$ is defined by a homogeneous 
% The sensitivity equations for the parameters are

\emph{Claim}: $J$ is a negative semi-definite matrix. TODO: show with data.


% appendix_agent.tex — \input{appendix_agent} in the appendix section

\section{Agent-in-the-loop parameter optimisation}
\label{app:agent}

\paragraph{Architecture.}
The optimisation space combining simulation parameters (noise $\sigma$, connectivity
constraints, $n_{\mathrm{frames}}$) and machine learning parameters
(e.g. for GNN, $\lambda_0,\lambda_1,\gamma_0,\gamma_1,\mu_0,\mu_1$, learning rates, etc.) is too
large for grid search. We instead deploy a closed loop coupling three components:
external GNN experiments, LLM-agent (Claude code), and explicit structured
summaries. The system operates through text file exchange. The
instruction file defines the inverse problem, the metrics to track, the parameter ranges, and logging format. The
structured summary maintains accumulated understanding in two tiers:
established principles and open questions that persist across blocks, and a
working record of current hypotheses and recent iterations. A configuration
file is edited by the agent to set the next experiment; a results log is written
by the experiment engine and read by the agent.

\paragraph{Exploration protocol.}
Experiments are organised into blocks of 12 iterations. Within a block,
the LLM fixes the simulation regime and explores GNN training parameters via
At block boundaries the LLM selects the next simulation
regime to address gaps in the search space and compresses the block into the
structured summary, adding confirmed principles and retiring falsified
hypotheses. The structured summary also serves as an explicit long-term memory,
overcoming the finite context window of the agent and enabling accumulation across hundreds of iterations.

\paragraph{Scientific reasoning constraint.}
The agent is explicitly instructed not to predict the outcome of a model
training. Only experimental results can validate or falsify hypotheses. This
constraint distinguishes the system from prompt-optimisation approaches such as
OPRO \cite{yang_large_2024}: the LLM's role is hypothesis generation and
interpretation, not numerical optimisation. The epistemic reasoning —
induction, abduction, deduction, falsification — emerges from the interaction
between the three components rather than residing solely in the LLM.

\paragraph{Illustrative reasoning cycle.}
Table~\ref{tab:agent_example} shows a representative hypothesis→experiment→
revision cycle extracted from the exploration of the flyvis regime.

\begin{table}[ht!]
\centering\small
\begin{tabular}{p{2.2cm}p{9.5cm}}
\toprule
\textbf{Mode} & \textbf{Event} \\
\midrule
Abduction     & Effective rank $r=10$ (Dale's law regime) limits $R^2_{\mathbf{W}}$
                ceiling to $\approx 0.91$ regardless of learning rates. \\
Deduction     & Increasing $n_{\mathrm{frames}}$ should raise $r$ and lift the
                ceiling: predict $R^2_{\mathbf{W}} > 0.95$ at
                $n_{\mathrm{frames}} = 20{,}000$. \\
Falsification & $n_{\mathrm{frames}} = 20{,}000$ raises $r$ from 10 to 27 in
                rare stochastic initialisations only; ceiling not reliably
                broken. \\
Refinement    & Hypothesis revised: $n_{\mathrm{frames}}$ boosts $r$ only when
                connectivity rank exceeds a threshold; Dale's law stochasticity
                limits reliability. \\
Causal chain  & $\texttt{connectivity\_rank} \to r \to R^2_{\text{ceiling}} \to
                \texttt{convergence}$ established at iteration 48. \\
\bottomrule
\end{tabular}
\caption{Representative abduction→deduction→falsification→refinement cycle.
  The causal chain linking effective rank to convergence persisted across all
  subsequent blocks as an established principle.}
\label{tab:agent_example}
\end{table}

\paragraph{Epistemic outcomes.}
Across 2,048 GNN optimisations spanning 128 neural activity regimes, the agent
accumulated the following ....



%%%%%%%%%%%%%%%%%%%%%%%%%%%%%%%%%%%%%%%%%%%%%%%%%%%%%%%%%%%%
\newpage
\input{checklist.tex}

\end{document}